\begin{frame}{Änderungsanträge (SEDE)}
\tiny
\vspace{-0.45cm}
\renewcommand{\arraystretch}{0.9}
\setlength{\tabcolsep}{2pt}
\setlength{\abovetopsep}{0pt}
\setlength{\belowrulesep}{0pt}
\begin{tabularx}{\textwidth}{@{}>{\centering\arraybackslash}m{1.55cm} >{\RaggedRight\arraybackslash}m{\dimexpr\linewidth-1.55cm-2\tabcolsep\relax}@{}}
\toprule
\multicolumn{1}{m{1.55cm}}{\centering\textbf{Fraktion}} & \multicolumn{1}{m{\dimexpr\linewidth-1.55cm-2\tabcolsep\relax}}{\centering\textbf{Änderungsantrag}} \\
\midrule
\includegraphics[width=2cm,height=1cm,keepaspectratio]{Bilder/EVP.png} & 1. Jeder EU-Mitgliedstaat soll sich verpflichten, einen Teil seiner nationalen Streitkräfte und Ausrüstung in eine gemeinsame Europäische Armee zu \textcolor{red}{\sout{übertragen, mit dem langfristigen Ziel der kompletten Auflösung aller nationaler Streitkräfte;}}\textcolor{green}{\uline{übertragen;}} ii) 2040 25 Prozent, iii) 2050 \textcolor{red}{\sout{50}}\textcolor{green}{\uline{40}} Prozent, \textcolor{red}{\sout{iv) 2070 100 Prozent}} seiner Streitkräfte und militärischen Güter in die Europäische Armee übertragen. a) Die Befehls- und Kommandogewalt \textcolor{green}{\uline{für die an die Europäische Armee übertragenen Teile }}obliegt dem Hohen Vertreter der EU für Außen- und Sicherheitspolitik als Oberbefehlshaber\textcolor{green}{\uline{. Die verbleibenden Teile bleiben unter nationaler Befehlsgewalt}}. \\
\includegraphics[width=2cm,height=1cm,keepaspectratio]{Bilder/SD.png} & werden muss, um im Verteidigungsfall eine effektive gemeinsame Verteidigung des Unionsgebiets zu ermöglichen;\textcolor{green}{\uline{ gleichzeitig sollen bereits bestehende NATO-Strukturen beibehalten und weiterhin gefördert werden.}} \\
\includegraphics[width=2cm,height=1cm,keepaspectratio]{Bilder/RE.png} & \textcolor{red}{\sout{a) Die Befehls- und Kommandogewalt obliegt dem Hohen Vertreter der EU für Außen- und Sicherheitspolitik als Oberbefehlshaber.}} b) Die Budgethoheit über die gemeinsamen Streitkräfte obliegt dem Europäischen Parlament. \textcolor{red}{\sout{Der Ausschuss für Sicherheit und Verteidigung (SEDE) soll mit der Kontrolle der Streitkräfte betraut werden.}}\textcolor{green}{\uline{Das EU-Parlament soll mit der Kontrolle der Streitkräfte betraut werden.Die Befehls- und Kommandogewalt obliegt einem Rat bestehend aus den Regierungschefs der EU-Mitglieder.}} \\
\includegraphics[width=2cm,height=1cm,keepaspectratio]{Bilder/PfE.png} & i) 2030 10 Prozent, ii) 2040 \textcolor{red}{\sout{25}}\textcolor{green}{\uline{20}} Prozent, iii) 2050 \textcolor{red}{\sout{50}}\textcolor{green}{\uline{30}} Prozent, iv) \textcolor{green}{\uline{2060 40 Prozent, v) }}2070 \textcolor{green}{\uline{50 Prozent, vi) 2080 55 Prozent, vii) 2090 60 Prozent, viii) 2100 65 Prozent, ix) 2110 70 Prozent, x) 2120 75 Prozent, xi) 2130 80 Prozent, xii) 2140 85 Prozent, xiii) 2150 90 Prozent, xiv) 2160 95 Prozent, xv) 2170 }}100 Prozent\textcolor{green}{\uline{, }} seiner Streitkräfte und militärischen Güter in die Europäische Armee übertragen. \\
\includegraphics[width=2cm,height=1cm,keepaspectratio]{Bilder/Left.png} & ii) 2040 25 Prozent, iii) 2050 50 Prozent, \textcolor{red}{\sout{iv) 2070 100 Prozent}} seiner Streitkräfte und militärischen Güter in die Europäische Armee übertragen. \\
\bottomrule
\end{tabularx}
\end{frame}

